% discussion.tex --- Section 7: Discussion
% ASSERT_CONVENTION: natural_units=kB_equals_1, entropy_base=nats, state_normalization=Tr(rho)=1, sequential_product=a&b=sqrt(a)b*sqrt(a), von_neumann_entropy=S(rho)=-Tr(rho*ln*rho), information=I(B;M)=S(B)+S(M)-S(BM), free_energy=F=E-TS

%----------------------------------------------------------------------
\section{Discussion}\label{sec:discussion}
%----------------------------------------------------------------------

%----------------------------------------------------------------------
\subsection{What this paper establishes and what remains open}
\label{sec:resolves}
%----------------------------------------------------------------------

The central results of this paper are the entropy gradient theorem
(Theorem~\ref{thm:entropy-gradient}), which constrains \emph{all}
self-modelers via three independent routes, and the structural
convergence of the thermodynamic and algebraic chains on
time-orientation, which narrows Gap~C: complexification is necessary
for Standard Model--like observers within $\hthree$, but the status of
non-SM self-modelers in non-complexified blocks remains open.

Two independent arguments lead to time-orientation.  The
\emph{thermodynamic chain} runs: self-modeling $\to$ Landauer bound
$\to$ free energy $\to$ entropy gradient $\to$ time-orientation
(Theorem~\ref{thm:entropy-gradient}).  The \emph{algebraic chain}
runs: complexification $\to$ chirality $\to$ time-orientation
(Theorem~\ref{thm:three-consequence}).  The convergence of these
chains on the same geometric structure---time-orientation---is a
structural insight connecting observer thermodynamics to particle
physics.  It is not a single deductive chain in which one implies the
other.

This convergence narrows Gap~C but does not close it.  For SM-like
observers---those whose matter content is described by the chiral
representation arising from $\Cl{6}$---complexification of $\Vhalf$ is
necessary, because without it the chirality operator has no
eigenstates and the $L/R$ decomposition does not exist.  For such
observers, the algebraic chain is blocked.  For the same observers,
the thermodynamic chain independently requires time-orientation,
which in the SM-like setting presupposes the geometric structures that
complexification provides.

The nine-link chain (Table~\ref{tab:chain}) is accordingly updated:
six links are proved algebraically (L1--L3, L5, L8, L9), one is
updated with a geometric consequence (L7), one is narrowed via
selection for SM-like observers (L6, this paper), and one remains
assumed (L4, Gap~B1).

We emphasize what remains open.  Four gaps are not addressed here:
\begin{itemize}
  \item \textbf{Gap~A} ($\hthree$ identification): Why is the
    exceptional Jordan algebra $\hthree$ the ``universe algebra''?
    Paper~5's non-composability argument~\cite{Paper5} motivates this
    choice but does not uniquely determine it.

  \item \textbf{Gap~B1} (rank-1 idempotent): What selects the
    observer's idempotent $E_{11}$?  All rank-1 idempotents are
    $\Ffour$-conjugate, so any selection principle must invoke
    additional structure beyond the algebra's automorphism group.

  \item \textbf{Gap~B2} (complex structure $u$): What selects the
    specific $u\in S^6$?  The moduli space $S^6$ is six-dimensional;
    the spectral action might generate a potential with a preferred
    minimum.

  \item \textbf{Gap~SA} (spectral action): The algebraic data
    established by Papers~5--8---gauge group, chiral representation,
    and complexified Peirce space---are the \emph{input} to the
    spectral action principle~\cite{CCM2007}, which should produce
    the full GR~+~SM Lagrangian as \emph{output}.  This computation
    has not been performed for the $\hthree$-based spectral triple.
\end{itemize}

\noindent
The selection narrowing of Gap~C is explicitly weaker than an
algebraic proof would be.  An algebraic proof would close Gap~C
unconditionally: $\Vhalf$ \emph{must} complexify as a consequence of
the algebra's structure.  The selection argument narrows it
conditionally: among SM-like blocks with $\rhoexp > 0$, $\Vhalf$ is
complexified.  This is a \emph{necessary condition} for SM-like
particle physics (complexification is required for chirality), not
a \emph{sufficient condition} (complexification alone does not
guarantee a viable self-modeler).

\paragraph{What this paper does not claim.}
We state four non-claims explicitly.
\begin{enumerate}
  \item This paper does \emph{not} prove that $\rhoexp = 0$ for all
    non-complexified blocks.  The entropy gradient theorem applies to
    all self-modelers regardless of their matter content; a
    non-complexified block could in principle support a non-SM
    self-modeler that satisfies the thermodynamic chain without chiral
    fermions.
  \item This paper does \emph{not} prove that all experience requires
    complexification.  The experiential measure $\rhoexp$ depends on
    the self-modeling axioms (L4, Gap~B1), and self-modelers with
    non-SM physics might exist under those axioms.
  \item Non-SM self-modelers---if they exist---are inaccessible to us:
    our observations are mediated by the SM-like physics that the
    algebraic chain describes.  They are not forbidden, merely outside
    the scope of this framework's predictions.
  \item The program explains \emph{our} physics, not all possible
    physics.  This has the same epistemic status as anthropic arguments
    for spatial dimensionality: the framework constrains the physics
    \emph{compatible with observers like us}, not the totality of
    possible physics.
\end{enumerate}

%----------------------------------------------------------------------
\subsection{Selection vs.\ algebraic derivation}
\label{sec:selection-vs-algebra}
%----------------------------------------------------------------------

The epistemic status of the Gap~C narrowing deserves careful
statement.  The experiential measure $\rhoexp$ selects complexified
blocks among SM-like observers, but does not force complexification
algebraically.  The algebraic obstruction---the $V_1 =
\mathbb{R}\cdot E_{11}$ bottleneck in the Peirce
decomposition~\cite{BGW2020}---remains genuine.  No algebraic
mechanism within the Jordan structure of $\hthree$ promotes the real
$\Vhalf = \mathbb{O}^2$ to the complex $\Vhalf^{\mathbb{C}} =
S_{10}^+$.  What the two-chain argument provides is a
\emph{structural} filter for SM-like observers: the algebraic chain
shows that such observers require complexification for chirality, and
the thermodynamic chain independently shows that all self-modelers
require time-orientation.  The convergence of both chains on the same
geometric prerequisite makes the case that SM-like self-modelers
inhabit complexified blocks.  Whether non-complexified blocks can
support non-SM self-modeling is a separate question that this paper
leaves open.

An analogy clarifies the logical structure.  In evolutionary biology,
natural selection does not force specific mutations to occur; it acts
on the phenotypic consequences of mutations that have already occurred.
Similarly, the experiential measure does not force structure space to
contain complexified blocks; it acts on the thermodynamic consequences
of blocks that are already present.  The selection is over
\emph{existing} structure, not a dynamical process that creates new
structure.

This connects to Hoffman's interface theory of
perception~\cite{Hoffman2015}.  In the self-modeling framework, a
block with $\rhoexp = 0$ is a passive interface: it exists in
structure space but supports no self-modeling activity.  A block with
$\rhoexp > 0$ is an active self-model: the body--model correlations
are maintained at thermodynamic cost, and the experiential measure
registers positive weight.  The distinction between passive and active
is thermodynamic, not fitness-theoretic: what separates $\rhoexp > 0$
from $\rhoexp = 0$ is the availability of an entropy gradient, not
an external fitness function.  For SM-like observers, the additional
algebraic requirement of chirality ties the thermodynamic condition
to complexification specifically.  For hypothetical non-SM
self-modelers, the thermodynamic condition still applies but the
algebraic link to complexification does not.

%----------------------------------------------------------------------
\subsection{The chirality--thermodynamics nexus}
\label{sec:nexus-discussion}
%----------------------------------------------------------------------

Corollary~\ref{cor:nexus} identifies time-orientation as the common
geometric prerequisite for both the chirality of the Standard Model
and the thermodynamic arrow of time.  This is a structural connection,
not a coincidence.

The chirality side of the connection is algebraic and geometric:
the choice $u\in S^6$ determines the internal chirality via the
$\Cl{6}$ volume form $\omegasix$~\cite{Paper7}, and its physical
realization as spacetime Weyl spinors requires a globally consistent
sign for the chirality operator $\Gammastar$, which in turn requires
time-orientation
(Proposition~\ref{prop:chirality-flip},
Theorem~\ref{thm:three-consequence}).
The thermodynamic side is dynamical: the entropy gradient theorem
(Theorem~\ref{thm:entropy-gradient}) shows that self-modelers require
$S(t) < \Smax$, defining a preferred temporal direction (the
``future'' is the direction of entropy increase), which on a
Lorentzian manifold is a time-orientation.

Both requirements trace back to self-modeling.  The algebraic choice
$u$ is made within the $\hthree$ framework that the self-modeling
axioms motivate~\cite{Paper5,Paper7}.  The entropy gradient arises
from the thermodynamic cost of self-modeling
(Theorem~\ref{thm:landauer}).  Time-orientation is the geometric
structure that both consequences demand, and it is this demand---not
temporal coincidence---that links particle physics chirality to the
arrow of time.

%----------------------------------------------------------------------
\subsection{CP violation: a structural remark}
\label{sec:cp-discussion}
%----------------------------------------------------------------------

The chirality--time-orientation link established in
Sec.~\ref{sec:chirality} admits a qualitative extension to CP
violation.  The CPT theorem (valid in any local Lorentz-invariant
quantum field theory) implies that CP violation entails T violation.
T violation---the dynamical asymmetry between forward and backward
temporal evolution---can only be defined on a time-oriented spacetime,
since ``forward'' and ``backward'' presuppose a preferred temporal
direction.  The resulting structural triangle is:
\begin{enumerate}
  \item Chirality $\leftrightarrow$ time-orientation
    (Theorem~\ref{thm:three-consequence});
  \item Self-modeling $\to$ entropy gradient $\to$ time-orientation
    (Theorem~\ref{thm:entropy-gradient});
  \item CP violation $\to$ T violation $\to$ time-orientation
    (CPT theorem).
\end{enumerate}
All three arrows converge on time-orientation as a shared geometric
prerequisite.  This suggests that chirality, the thermodynamic arrow
of time, and CP violation are structurally connected through their
common dependence on the same spacetime structure.

We state explicitly what this structural remark does \emph{not} claim.
The framework does not derive the magnitude of CP violation: the
Jarlskog invariant, the CKM phase, and the strong CP parameter are
Standard Model quantities determined by Yukawa couplings, which lie
outside the scope of the self-modeling axioms.  The framework does not
derive baryon asymmetry from the entropy gradient: Sakharov's
conditions~(a) (baryon number violation) and~(b) (CP violation) are
independent requirements that the entropy gradient theorem does not
address.  And no equation in this framework relates the entropy
gradient rate to the CP violation magnitude.  The connection is
structural and qualitative: all three phenomena require
time-orientation, and the self-modeling framework shows why
time-orientation is necessary.

%----------------------------------------------------------------------
\subsection{Assumption hierarchy}\label{sec:assumptions-discussion}
%----------------------------------------------------------------------

The master theorem (Theorem~\ref{thm:master}) rests on seven
assumptions (Table~\ref{tab:assumptions}), and the quantitative
predictions of Sec.~\ref{sec:predictions} require three additional
cosmological inputs (A8--A10).  The assumptions are not equally
strong; their failure modes differ qualitatively.

Assumption~A3 (closed system equilibration) is the strongest.  It
connects local thermodynamics---the Landauer cost per self-modeling
cycle---to a cosmological conclusion---the necessity of a low-entropy
past.  If A3 fails because an infinite free energy source exists
(e.g., an infinite-dimensional Hilbert space, or unknown physics that
provides work without depleting a reservoir), the chain breaks at
Link~3 of Proposition~\ref{prop:chain}: the entropy gradient theorem
still holds via Routes~B and~C, but Route~A's quantitative
connection to the Landauer cost is weakened.

Assumptions~A6--A7 (continuum limit and Lorentzian signature) are
standard in the sense that they are shared by essentially all
approaches to emergent spacetime, but they remain unvalidated from
first principles within the self-modeling framework.  Paper~6's
lattice~\cite{Paper6} provides the microscopic structure; the claim
that its continuum limit yields a smooth Lorentzian manifold is
physically motivated by Lieb-Robinson bounds but not rigorously
proved.  If A6 or A7 fail, the three-consequence theorem
(Theorem~\ref{thm:three-consequence}) reduces to Paper~7's
two-consequence theorem: consequences~(a) and~(b) (gauge group and
chirality) survive as purely algebraic results, while consequence~(c)
(time-orientation requirement) and the complexification selection
argument are lost.

Assumptions~A8--A10 (cosmological entropy budget, Penrose's initial
entropy estimate, and self-modeling parameters) are external inputs to
the framework, not consequences of it.  They enter only the
order-of-magnitude estimates P7--P9 of
Sec.~\ref{sec:predictions} and do not affect the framework
predictions P1--P6.  Their uncertainty is large (several orders of
magnitude in each parameter), but the qualitative conclusions are
robust: the 94-order-of-magnitude gap between the Landauer deficit and
the Penrose deficit (Eq.~\eqref{eq:gap-94}) survives any plausible
variation of these inputs.

%----------------------------------------------------------------------
\subsection{Future directions}\label{sec:future}
%----------------------------------------------------------------------

Five concrete extensions of the present work are identified, each
connected to a specific open problem.

\medskip\noindent
\textbf{1. Spectral action computation (Gap~SA).}
The algebraic data established by Papers~5--8 provide the input to
the spectral action principle~\cite{CCM2007}: the gauge group
$\SMgauge$, the chiral representation from $\Cl{6}$, and the
complexified Peirce space $S_{10}^+$.  Computing the spectral action
for an $\hthree$-based spectral triple would yield the full GR~+~SM
Lagrangian, including the Higgs sector, Yukawa couplings, and the
cosmological constant.  The principal challenge is the
non-associativity of the octonions, which prevents a straightforward
application of the Connes--Chamseddine--Marcolli
formalism~\cite{CCM2007}.

\medskip\noindent
\textbf{2. Sharper Landauer bound from detailed self-modeling models.}
The Landauer bound of Theorem~\ref{thm:landauer} gives the minimum
thermodynamic cost of self-modeling per update cycle.  This bound
is tight for quasi-static erasure protocols but may be far from tight
for realistic self-modeling dynamics.  A detailed model of the
self-modeling cycle---specifying the measurement, erasure, and writing
steps concretely---would yield a tighter bound, potentially shifting
the 94-order-of-magnitude gap of Eq.~\eqref{eq:gap-94}.

\medskip\noindent
\textbf{3. Generation structure.}
The present framework produces one generation of chiral fermions in
the $\rep{16}$ of $\Spin{10}$.  The Standard Model requires three
generations.  The ``3'' in $h_3(\mathbb{O})$ (the algebra of
$3\times 3$ matrices) is suggestive but no mechanism producing three
copies of the $\rep{16}$ has been established.
Furey~\cite{Furey2018} proposed a route via the full octonion algebra
acting on $\Cl{6}$; Boyle~\cite{Boyle2020} explored triality.
Making any of these connections rigorous within the self-modeling
framework is an open challenge.

\medskip\noindent
\textbf{4. Quantitative CP--entropy connection.}
The structural triangle of Sec.~\ref{sec:cp-discussion} identifies a
qualitative link between CP violation and the entropy gradient.  A
quantitative version would require showing that gauge groups
admitting CP-violating phases are selected for by the experiential
measure---an argument analogous to the Gap~C selection but operating
at the level of the CKM structure rather than the complexification
step.  This lies well beyond the scope of the present work.

\medskip\noindent
\textbf{5. Experimental and observational signatures.}
The predictions of Sec.~\ref{sec:predictions} are primarily
consistency checks rather than novel observational predictions.  The
most direct path to testable consequences is through the spectral
action computation (direction~1), which would predict specific
relations among Standard Model parameters.  A less direct but
potentially more distinctive signature is the exhaustion timescale
$\tauex$: any observation constraining the long-term thermodynamic
fate of the universe constrains the self-modeling capacity.

%----------------------------------------------------------------------
\subsection{Concluding remarks}\label{sec:concluding}
%----------------------------------------------------------------------

The question ``why complex numbers in quantum mechanics?'' has
been posed since the formalism's inception.  Within the self-modeling
framework, this paper provides a partial answer for SM-like observers:
because chirality requires complexification, and chirality requires the
same time-orientation that self-modeling independently demands.

Two independent chains converge on time-orientation.  The
thermodynamic chain: self-modeling demands an entropy gradient
(Theorem~\ref{thm:entropy-gradient}), which defines a
time-orientation on the emergent spacetime.  The algebraic chain:
the choice $u \in S^6$ determines the gauge group and chirality
(Theorem~\ref{thm:three-consequence}), and the physical realization
of chirality as spacetime Weyl spinors requires time-orientation
(Proposition~\ref{prop:chirality-flip}).  The complexification of
$\Vhalf$ enters through the algebraic chain: without it, the $L/R$
decomposition from $\Cl{6}$ does not exist, and the chirality that
requires time-orientation is absent.

The complexification step that seemed algebraically unmotivated in
Paper~7~\cite{Paper7} is thereby given a physical rationale: it is
necessary for SM-like observers, whose chiral matter content requires
the geometric structure that self-modeling thermodynamics independently
demands.  This is not logical closure---non-SM self-modelers in
non-complexified blocks are not excluded---but it is structural
convergence: two independent arguments, one thermodynamic and one
algebraic, identify time-orientation as indispensable, and for
SM-like physics, time-orientation requires complexification.
